\chapter{مقدمه}
% sec 1
\section{بیان مسئله و اهمیت آن}
تخمین عمق صحنه از روی یک تصویر تک‌دوربینه یکی از مسائل بنیادی و در عین حال چالش‌برانگیز در بینایی ماشین\enfootnote{Computer Vision} به شمار می‌رود~\cite{Saxena2009Make3D,Eigen2014Depth}.
این مسئله در کاربردهای متنوعی از جمله رباتیک متحرک\enfootnote{Mobile Robotics}، ناوبری خودکار\enfootnote{Autonomous Navigation}، بازسازی سه‌بعدی\enfootnote{3D Reconstruction}، واقعیت افزوده\enfootnote{Augmented Reality} و سیستم‌های کمک‌راننده\enfootnote{Advanced Driver-Assistance Systems (ADAS)} نقش کلیدی ایفا می‌کند~\cite{Godard2017Unsupervised,Ranftl2019Robust}.
برخلاف روش‌های مبتنی بر چند نما یا حسگرهای فعال،
تخمین عمق تک‌دوربینه\enfootnote{Monocular Depth Estimation} تنها بر اطلاعات یک تصویر دوبعدی متکی است و از این‌رو ذاتاً با ابهام‌های هندسی متعددی مواجه می‌شود؛ یک تصویر منفرد تنها پرتوِ نگاه هر پیکسل را معین می‌کند و نه فاصلهٔ آن را~\cite{HartleyZisserman2004,Eigen2014Depth}.

با پیشرفت روش‌های یادگیری عمیق، به‌ویژه شبکه‌های عصبی کانولوشنی\enfootnote{Convolutional Neural Networks} و مدل‌های مبتنی بر ترنسفورمر،
کیفیت تخمین عمق تک‌دوربینه به‌طور قابل توجهی بهبود یافته است~\cite{Eigen2015Predicting,Ranftl2021DPT}.
بسیاری از مدل‌های مدرن قادرند ساختار نسبی صحنه را با دقت بالا بازسازی کنند و روابط عمقی میان اجسام را به‌خوبی بیاموزند. با این حال، خروجی این مدل‌ها در اغلب موارد فاقد مقیاس فیزیکی واقعی بوده و به‌صورت \emph{عمق نسبی} یا \emph{\lr{non-metric depth}} بیان می‌شود. این محدودیت باعث می‌شود که استفادهٔ مستقیم از چنین مدل‌هایی در کاربردهایی که نیازمند اندازه‌گیری‌های متریک دقیق هستند، با چالش مواجه شود.

مسئلهٔ اصلی در این میان، ابهام مقیاس ذاتی در تخمین عمق تک‌دوربینه است.
در غیاب اطلاعات هندسی صریح دربارهٔ دوربین تصویربرداری،
بازیابی عمق متریک تنها بر اساس شدت پیکسل‌ها ذاتاً نامعین است~\cite{HartleyZisserman2004}. اگرچه برخی روش‌ها با اعمال نرمال‌سازی یا استفاده از معیارهای نسبی، عملکرد مناسبی در سنجش‌های غیرمتریک نشان می‌دهند، اما این رویکردها تضمینی برای پایداری مقیاس فیزیکی عمق فراهم نمی‌کنند. در نتیجه، تغییر شرایط تصویربرداری، از جمله تفاوت در نوع دوربین یا پارامترهای آن، می‌تواند منجر به تغییرات معنادار در کیفیت تخمین عمق شود.


\begin{figure}[htbp]
    \centering
    \begin{tikzpicture}[font=\small]

    %%% ---------- پنل (الف): ابهام مقیاس روی یک پرتو دید (سمت راست) ----------
    \begin{scope}[xshift=7.7cm]
        \fill[thPanel, rounded corners=3pt] (-0.55,-2.35) rectangle (6.35,2.45);

        % بدنهٔ دوربین و مرکز نوری
        \draw[thick, draw=thInk, fill=white] (-0.10,-0.30) rectangle (0.50,0.30);
        \fill[thInk] (0.50,0) circle (1.6pt);
        \node[thnote, below=2pt] at (0.20,-0.32) {$C$};

        % صفحهٔ تصویر
        \draw[line width=1.4pt, thProposed] (1.35,-0.95) -- (1.35,0.95);
        \node[thnote, text=thProposed, anchor=south] at (1.35,1.00) {\rl{صفحهٔ تصویر}};

        % پرتوهای حاشیهٔ میدان دید
        \draw[dashed, thInk!65] (0.50,0) -- (5.90,1.45);
        \draw[dashed, thInk!65] (0.50,0) -- (5.90,-1.45);

        % تصویر مشترک هر دو جسم روی صفحهٔ تصویر
        \draw[line width=2pt, thAccent] (1.35,-0.228) -- (1.35,0.228);

        % جسم نزدیک و کوچک
        \fill[thControl, opacity=0.9] (2.28,-0.51) rectangle (2.52,0.51);
        \node[thnote, text=thControl, anchor=north, align=center] at (2.40,-0.60)
            {\rl{جسم نزدیک}\\ \rl{و کوچک}};

        % جسم دور و بزرگ
        \fill[thIdeal, opacity=0.9] (4.88,-1.21) rectangle (5.12,1.21);
        \node[thnote, text=thIdeal, anchor=north, align=center] at (5.00,-1.30)
            {\rl{جسم دور}\\ \rl{و بزرگ}};

        \node[font=\small, anchor=south] at (2.90,1.95) {\rl{(الف) ابهام مقیاس روی یک پرتو دید}};
    \end{scope}

    %%% ---------- پنل (ب): هم‌ارزی $(f,Z)$ و $(\alpha f, \alpha Z)$ (سمت چپ) ----------
    \begin{scope}
        \fill[thPanelAlt, rounded corners=3pt] (-0.55,-2.35) rectangle (6.35,2.45);

        \node[thblock, fill=white, text width=2.55cm] (camA) at (1.15,1.05)
            {\rl{دوربین ۱}\\[1pt] \rl{کانون $f$}\\ \rl{عمق صحنه $Z$}};
        \node[thblock, fill=white, text width=2.55cm] (camB) at (1.15,-1.05)
            {\rl{دوربین ۲}\\[1pt] \rl{کانون $\alpha f$}\\ \rl{عمق صحنه $\alpha Z$}};

        \node[thdata, fill=white, text width=1.9cm] (img) at (4.75,0)
            {\rl{تصویر ثبت‌شدهٔ یکسان}};

        \draw[tharrow] (camA) -- (img);
        \draw[tharrow] (camB) -- (img);

        \node[thnote, anchor=north, text width=5.7cm] at (2.90,-1.75)
            {$u=\dfrac{f X}{Z}=\dfrac{\alpha f X}{\alpha Z}$};

        \node[font=\small, anchor=south] at (2.90,1.95) {\rl{(ب) هم‌ارزی دقیق کانون و مقیاس}};
    \end{scope}

    \end{tikzpicture}
    \caption[ماهیت هندسی ابهام مقیاس در تخمین عمق تک‌دوربینه]{ماهیت
    هندسی ابهام مقیاس در تخمین عمق تک‌دوربینه.
    (الف) دو جسم با اندازه و فاصلهٔ متفاوت که روی یک پرتو دید قرار دارند،
    تصویر کاملاً یکسانی تولید می‌کنند؛ بنابراین شدت پیکسل‌ها به‌تنهایی
    مقیاس مطلق صحنه را تعیین نمی‌کند.
    (ب) صورت کمی همین ابهام: اگر فاصلهٔ کانونی و تمام عمق‌های صحنه
    هم‌زمان در ضریب $\alpha$ ضرب شوند، تصویر حاصل بدون هیچ تغییری باقی می‌ماند.
    رفع این ابهام تنها با دانستن فاصلهٔ کانونی --- یعنی بخشی از ماتریس
    مشخصهٔ دوربین $K$ --- ممکن است. همین هم‌ارزی، در فصل سوم مبنای
    افزونه‌سازی «لرزش کانونی» قرار می‌گیرد.}
    \label{fig:scale_ambiguity}
\end{figure}

از این منظر، دستیابی به تخمین عمق متریک\enfootnote{Metric Depth} پایدار و قابل تعمیم، همچنان یکی از چالش‌های باز در حوزهٔ بینایی ماشین محسوب می‌شود. پرداختن به این مسئله نه‌تنها از دیدگاه نظری اهمیت دارد، بلکه از منظر عملی نیز برای توسعهٔ سیستم‌های قابل اعتماد در محیط‌های واقعی ضروری است. این پژوهش با تمرکز بر این چالش اساسی، به بررسی راهکارهایی برای ارتقای مدل‌های تخمین عمق تک‌دوربینه از خروجی‌های غیرمتریک به عمق متریک پایدار می‌پردازد.

% sec 2
\section{چالش‌های عمق متریک در مدل‌های یادگیری عمیق}
اگرچه مدل‌های یادگیری عمیق در سال‌های اخیر پیشرفت چشمگیری در تخمین عمق تک‌دوربینه داشته‌اند،
اما اغلب این مدل‌ها به‌طور ضمنی برای یادگیری ساختار نسبی صحنه طراحی شده‌اند~\cite{Eigen2014Depth,Ranftl2019Robust}. بسیاری از معیارهای ارزیابی رایج نیز بر سنجش خطاهای نرمال‌شده یا نسبی تمرکز دارند
و به‌جای مقیاس فیزیکی مطلق، همبستگی ساختاری عمق را مورد توجه قرار می‌دهند~\cite{Ranftl2019Robust,Ranftl2021DPT}. این موضوع سبب شده است که بهبود عملکرد در این معیارها لزوماً به معنای دستیابی به عمق متریک دقیق نباشد.

یکی از چالش‌های مهم در این زمینه، وابستگی مدل‌ها به توزیع داده‌های آموزشی است. مدل‌هایی که بر روی یک یا چند دیتاست خاص آموزش داده می‌شوند، معمولاً فرض ضمنی ثبات شرایط تصویربرداری را در خود دارند. در عمل، تغییر نوع دوربین، فاصلهٔ کانونی یا میدان دید\enfootnote{Field of View (FoV)} می‌تواند توزیع نگاشت تصویر به عمق را تغییر دهد
و باعث افت عملکرد مدل شود~\cite{Facil2019CAMConvs,Yin2023Metric3D}. این مسئله به‌ویژه در معیارهای متریک که نسبت به مقیاس فیزیکی حساس هستند، نمود بیشتری پیدا می‌کند.

چالش دیگر، فقدان استفادهٔ صریح از اطلاعات هندسی دوربین در بسیاری از مدل‌های موجود است~\cite{Facil2019CAMConvs,Yin2023Metric3D}. در حالی که فرآیند تصویربرداری به‌طور مستقیم تحت تأثیر پارامترهای درونی دوربین\enfootnote{Camera Intrinsics} قرار دارد، اغلب روش‌های یادگیری عمیق این اطلاعات را نادیده می‌گیرند یا به‌صورت ضمنی در داده‌ها مستتر فرض می‌کنند. این رویکرد می‌تواند منجر به یادگیری روابطی شود که تنها در شرایط خاص معتبر بوده و قابلیت تعمیم به سناریوهای متفاوت را ندارند.

در مجموع، این چالش‌ها نشان می‌دهند که برای دستیابی به تخمین عمق متریک پایدار، صرف بهبود معماری شبکه یا افزایش حجم داده‌های آموزشی کافی نیست. بلکه لازم است به‌صورت آگاهانه به نقش هندسهٔ تصویربرداری و اطلاعات دوربین در فرآیند تخمین عمق توجه شود. این ضرورت، انگیزهٔ اصلی پژوهش حاضر را شکل می‌دهد و زمینه را برای ارائهٔ رویکردی فراهم می‌کند که بتواند پایداری عمق متریک را در مواجهه با تغییرات شرایط تصویربرداری بهبود بخشد.

% sec 3
\section{خلأ پژوهشی}
بخش عمده‌ای از پژوهش‌های پیشین بر بهبود دقت تخمین عمق نسبی تمرکز داشته‌اند~\cite{Eigen2014Depth,Godard2017Unsupervised} و موفق شده‌اند ساختار کلی صحنه و روابط عمقی میان اجسام را با کیفیت مناسبی بازسازی کنند. با این حال، در بسیاری از این مطالعات، مسئلهٔ مقیاس فیزیکی عمق یا به‌صورت ضمنی نادیده گرفته شده و یا به‌عنوان یک چالش ثانویه تلقی شده است.

از سوی دیگر، روش‌هایی که به تخمین عمق متریک پرداخته‌اند،
غالباً متکی بر فرض‌های محدودکننده‌ای نظیر ثبات نوع دوربین یا شرایط تصویربرداری کنترل‌شده~\cite{Ranftl2021DPT,Bhat2023ZoeDepth} یا دسترسی به داده‌های کالیبره‌شدهٔ خاص بوده‌اند. این فرض‌ها موجب می‌شود که قابلیت تعمیم چنین مدل‌هایی به سناریوهای واقعی و متنوع کاهش یابد. به‌ویژه، تأثیر تغییر پارامترهای درونی دوربین بر عملکرد مدل‌ها در بسیاری از این کارها به‌صورت نظام‌مند مورد مطالعه قرار نگرفته است.

نکتهٔ قابل توجه دیگر آن است که در اغلب پژوهش‌های موجود، 
ارزیابی مدل‌ها عمدتاً بر اساس معیارهای غیرمتریک\enfootnote{Non-metric Metrics} یا نرمال‌شده انجام شده است~\cite{Eigen2015Predicting}. این معیارها اگرچه برای سنجش کیفیت نسبی عمق مناسب هستند، اما قادر به آشکارسازی ناپایداری‌های مرتبط با مقیاس فیزیکی و تغییر شرایط تصویربرداری نیستند. در نتیجه، ممکن است مدلی از نظر این معیارها عملکرد مطلوبی داشته باشد، در حالی که در کاربردهای متریک و واقعی دچار خطای قابل توجه شود.

علاوه بر این، اگرچه نقش هندسهٔ دوربین در فرآیند تصویربرداری از دیرباز در هندسهٔ
چندنمایی صورت‌بندی شده است~\cite{HartleyZisserman2004}، در بدنهٔ اصلی مدل‌های
یادگیری عمیق تخمین عمق، این اطلاعات یا کاملاً نادیده گرفته می‌شوند یا فرض می‌شود
که شبکه می‌تواند اثر آن‌ها را به‌طور ضمنی از داده‌ها بیاموزد~\cite{Eigen2014Depth,Godard2019Digging,Ranftl2021DPT}.
در برابر این جریان غالب، خطی از پژوهش‌ها پارامترهای درونی دوربین را به‌صورت صریح
وارد مدل کرده‌اند: \lr{CAM-Convs} این پارامترها را به‌شکل کانال‌های مختصاتی به
کانولوشن‌ها می‌افزاید~\cite{Facil2019CAMConvs}، \lr{Metric3D} تصویر یا برچسب عمق
را به یک «فضای دوربین متعارف» نگاشت می‌کند~\cite{Yin2023Metric3D}، برخی روش‌ها مدل
انتشاری را بر میدان دید شرطی می‌کنند~\cite{Saxena2023FOV}، و \lr{UniDepth} یک
شاخهٔ دوربین را در کنار شاخهٔ عمق آموزش می‌دهد~\cite{Piccinelli2024UniDepth}.
بنابراین ادعای این پژوهش آن نیست که استفاده از $K$ ایدهٔ تازه‌ای است؛ خلأ در جای
دیگری است: در این کارها، ورود $K$ همواره با تغییر هم‌زمان معماری، مجموعهٔ آموزشی و
برنامهٔ آموزش همراه بوده و بنابراین \emph{سهم علّی خودِ $K$} از سایر عوامل تفکیک
نشده است.

بر این اساس، خلأ مورد هدف این پژوهش را می‌توان به‌صورت مشخص چنین بیان کرد: یک
\emph{مطالعهٔ کنترل‌شدهٔ زوجی} که در آن هستهٔ یک مدل بنیادین ثابت و فریزشده نگه
داشته شود، هر دو بازوی آزمایش دقیقاً داده، برنامهٔ آموزشی و بذر تصادفی یکسانی
دریافت کنند، و تنها متغیر مورد دستکاری، دسترسی شبکه به ماتریس مشخصهٔ دوربین باشد.
چنین طرحی امکان می‌دهد سهم $K$ با معیارهای متریک و بدون تراز مقیاس اندازه‌گیری
شود و با آزمون‌های حذف و تحریف $K$ راستی‌آزمایی گردد. پژوهش حاضر با تمرکز بر همین
خلأ، چارچوبی برای تحلیل و اندازه‌گیری سهم اطلاعات هندسی دوربین در پایداری عمق
متریک ارائه می‌دهد.

% sec 4
\section{اهداف و سوالات تحقیق}
با توجه به خلأهای شناسایی‌شده در ادبیات موضوع و چالش‌های موجود در تخمین عمق تک‌دوربینه، پژوهش حاضر با هدف بررسی امکان دستیابی به تخمین عمق متریک پایدار در چارچوب یادگیری عمیق انجام شده است. تمرکز اصلی این تحقیق بر تحلیل نقش هندسهٔ دوربین در فرآیند تخمین عمق و ارزیابی تأثیر لحاظ‌کردن اطلاعات درونی دوربین بر پایداری و دقت خروجی مدل‌ها است.

هدف کلی این پژوهش، بررسی این مسئله است که آیا می‌توان با استفاده از اطلاعات صریح مربوط به مشخصات درونی دوربین، محدودیت‌های ذاتی تخمین عمق غیرمتریک را کاهش داد و به تخمینی دست یافت که نسبت به تغییر شرایط تصویربرداری پایدارتر باشد. در راستای این هدف کلی، اهداف جزئی زیر دنبال می‌شوند:
\begin{itemize}
	\item
	بررسی میزان وابستگی مدل‌های تخمین عمق تک‌دوربینه به نوع دیتاست و پارامترهای دوربین؛
	\item
	تحلیل تفاوت رفتار معیارهای ارزیابی غیرمتریک و متریک در مواجهه با تغییرات شرایط تصویربرداری؛
	\item
	ارزیابی امکان بهبود دقت و پایداری تخمین عمق متریک با استفاده از اطلاعات هندسی صریح.
\end{itemize}

بر این اساس، سوالات اصلی تحقیق به‌صورت زیر مطرح می‌شوند:
\begin{enumerate}
	\item
	آیا مدل‌های رایج تخمین عمق تک‌دوربینه در برابر تغییر دیتاست و پارامترهای درونی دوربین رفتار پایداری از خود نشان می‌دهند؟
	\item
	آیا استفاده از معیارهای غیرمتریک می‌تواند ناپایداری‌های مرتبط با مقیاس فیزیکی عمق را پنهان کند؟
	\item
	آیا لحاظ‌کردن اطلاعات درونی دوربین می‌تواند به بهبود تخمین عمق متریک و کاهش حساسیت مدل نسبت به تغییرات دوربین منجر شود؟
\end{enumerate}

جدول~\ref{tab:rq_hypothesis_map} نگاشت این سوالات را به فرضیه‌های تجربی متناظر و به بخشی از پایان‌نامه که شاهد آن ارائه می‌شود نشان می‌دهد. این نگاشت، مسیر خواندن فصل چهارم را نیز مشخص می‌کند.

\begin{table}[htbp]
\centering
\caption[نگاشت سوالات تحقیق به فرضیه‌ها و شواهد]{نگاشت سوالات تحقیق به
فرضیه‌های تجربی متناظر و بخش‌هایی که شاهد تجربی هر یک در آن‌ها ارائه می‌شود.}
\label{tab:rq_hypothesis_map}
\small
\renewcommand{\arraystretch}{1.3}
\begin{tabularx}{\textwidth}{@{}c
  >{\raggedleft\arraybackslash}X
  >{\raggedleft\arraybackslash}X
  >{\raggedleft\arraybackslash}p{0.235\textwidth}@{}}
\toprule
ردیف & سوال تحقیق & فرضیهٔ تجربی متناظر & شاهد اصلی در این پایان‌نامه \\ \midrule
۱ &
پایداری مدل‌های رایج در برابر تغییر دیتاست و پارامترهای درونی دوربین &
فرضیهٔ اول: \basemodel{} در معیارهای متریک به دیتاست و دوربین حساس است &
جدول~\ref{tab:middlebury_camera_audit} و بخش~\ref{subsec:basic_model_cross_dataset} \\
۲ &
توانایی معیارهای غیرمتریک در پنهان‌کردن ناپایداری مقیاس فیزیکی &
فرضیهٔ سوم: بهبود متریک بدون افت ساختار نسبی به دست می‌آید و پروتکل نسبی نسبت به مقیاس نابیناست &
بخش‌های~\ref{subsec:aug_basic_results} و~\ref{subsec:protocol_capacity_dependence} \\
۳ &
اثر لحاظ صریح ماتریس مشخصهٔ دوربین بر دقت و پایداری عمق متریک &
فرضیهٔ دوم: ورودی صریح $K$ خطای متریک را به‌صورت معنادار کاهش می‌دهد &
بخش‌های~\ref{subsec:metric_protocol_results}، \ref{subsec:focal_jitter_results} و~\ref{subsec:focal_ablation} \\
۴ &
اثبات اینکه بهبود مشاهده‌شده واقعاً از مصرف $K$ می‌آید، نه هم‌بستگی جانبی &
شرط لازم برای اعتبار فرضیهٔ دوم: پاسخ خروجی به $K$ باید کمی و مطابق هندسهٔ روزنه‌ای باشد &
آزمون جاروی $K$ (\ref{subsec:intrinsics_variation}) و آزمون حذف $K$ (\ref{subsec:k_knockout}) \\
\bottomrule
\end{tabularx}
\end{table}

پاسخ به این سوالات، چارچوب اصلی تحلیل‌های تجربی ارائه‌شده در فصل چهارم را تشکیل می‌دهد و مبنای نتیجه‌گیری‌های نهایی پژوهش حاضر قرار می‌گیرد.

% sec 5
\section{ساختار پایان‌نامه}
این پایان‌نامه در پنج فصل سازمان‌دهی شده است که هر فصل به بخش مشخصی از فرآیند پژوهش اختصاص دارد.

در فصل اول، مقدمه‌ای بر مسئلهٔ تخمین عمق تک‌دوربینه ارائه شد و با تبیین اهمیت موضوع، چالش‌های موجود و خلأهای پژوهشی، اهداف و سوالات تحقیق مطرح گردید.

فصل دوم به مرور ادبیات موضوع اختصاص دارد. در این فصل، پژوهش‌های پیشین در حوزهٔ تخمین عمق تک‌دوربینه، روش‌های متریک و غیرمتریک، نقش هندسهٔ دوربین و معیارهای ارزیابی مورد بررسی و تحلیل قرار می‌گیرند و جایگاه پژوهش حاضر نسبت به کارهای موجود مشخص می‌شود.

در فصل سوم، روش پیشنهادی تحقیق به‌صورت جامع تشریح می‌شود. این فصل شامل صورت‌بندی مسئله، تحلیل هندسی فرآیند تصویربرداری، معرفی معماری پایه، انگیزه و چارچوب روش پیشنهادی، راهبرد آموزش، توابع هزینه و جزئیات پیاده‌سازی است.

فصل چهارم به ارائهٔ تنظیمات آزمایشی، نتایج کمی و کیفی و تحلیل آن‌ها اختصاص دارد. در این فصل، عملکرد مدل‌ها بر روی دیتاست‌های مختلف بررسی شده و پایداری تخمین عمق نسبت به تغییر پارامترهای دوربین به‌صورت تجربی تحلیل می‌شود.

شکل~\ref{fig:thesis_roadmap} این ساختار را به‌صورت یک نمای کلی نشان می‌دهد.

\begin{figure}[htbp]
    \centering
    \begin{tikzpicture}[font=\scriptsize, node distance=0pt]
    \tikzset{
      chap/.style={thblock, text width=2.30cm, minimum height=1.85cm, font=\scriptsize},
    }
    \node[chap, fill=thPanelGray] (c5) at (0.00,0) {\rl{\textbf{فصل ۵}}\\[2pt] \rl{بحث و نتیجه‌گیری}\\[2pt] \rl{نوآوری‌ها، محدودیت‌ها و کار آینده}};
    \node[chap, fill=thPanelGreen] (c4) at (2.85,0) {\rl{\textbf{فصل ۴}}\\[2pt] \rl{نتایج}\\[2pt] \rl{مطالعهٔ کنترل‌شده و آزمون‌های علّی}};
    \node[chap, fill=thPanel] (c3) at (5.70,0) {\rl{\textbf{فصل ۳}}\\[2pt] \rl{روش تحقیق}\\[2pt] \rl{هندسهٔ دوربین و ادغام صریح $K$}};
    \node[chap, fill=thPanelAlt] (c2) at (8.55,0) {\rl{\textbf{فصل ۲}}\\[2pt] \rl{مروری بر مطالعات}\\[2pt] \rl{پیشینه و خلأ پژوهشی}};
    \node[chap, fill=thPanelRed] (c1) at (11.40,0) {\rl{\textbf{فصل ۱}}\\[2pt] \rl{مقدمه}\\[2pt] \rl{مسئله، اهداف و سوالات}};

    \draw[tharrow] (c1) -- (c2);
    \draw[tharrow] (c2) -- (c3);
    \draw[tharrow] (c3) -- (c4);
    \draw[tharrow] (c4) -- (c5);

    \draw[thdashed, rounded corners=3pt]
        (c1.south) -- ++(0,-0.75) -| (c4.south);
    \node[thnote, text=thAccent, anchor=north] at (6.9,-0.80)
        {\rl{سوالات فصل نخست، مستقیماً در فصل چهارم آزمون می‌شوند}};
    \end{tikzpicture}
    \caption[ساختار پایان‌نامه]{ساختار پایان‌نامه و رابطهٔ منطقی میان فصل‌ها.
    سوالات و فرضیه‌های مطرح‌شده در فصل نخست (جدول~\ref{tab:rq_hypothesis_map})
    در فصل چهارم به‌صورت تجربی آزمون می‌شوند و فصل پنجم برداشت نهایی را ارائه می‌کند.}
    \label{fig:thesis_roadmap}
\end{figure}

در نهایت، فصل پنجم به بحث و نتیجه‌گیری اختصاص دارد. در این فصل، یافته‌های اصلی پژوهش جمع‌بندی شده، نوآوری‌های نظری و عملی تحقیق بیان می‌شود و پیشنهادهایی برای پژوهش‌های آتی ارائه می‌گردد.

